\paragraph{Reconcile financing movements and show the dates still covered.}
A verified company trace can explain the financing balances at the same time as it explains the operating cycle. `DeclaredFinancingSchedules/1` first recomputes the v2 program and chooses a successful report. It then rebuilds the declared financing definitions from the exact policy, intervals and role map. Registry names do not classify an event; the method definition hash does.

The schedule records debt issues and repayments, interest accruals and payments, capital contributions, dividend declarations and dividend payments. Each row retains its event date, typed inputs, method, request and result hashes. Signed ledger credits are presented as positive liabilities. Opening plus increases less decreases equals the selected closing balance for debt, interest payable, capital, retained direct movements and dividends.

Interest coverage is a separate result. In the example, the cursor starts on 1 January 2027 and advances through intervals of 90, 91 and 184 days to 1 January 2028. The three rows therefore cover 365 days. A balanced report with a cursor that stops earlier remains a valid provisional balance report with an explicit accrual gap. A carry cannot silently complete the missing suffix. Unknown methods that change a scoped account, tracked debt or cursor refuse classification, even when their net journal is zero.

The four operating schedules remain compatible with the v2 program. Their partial cash bridge leaves the unpaid-interest residual visible under the selected operating-interest policy. This is an explanation for the synthetic example, not a general indirect cash-flow rule. Contract maturity, fees, valuation, outstanding shares, lawful distributions, source completeness and full notes remain unresolved.
